A brain--computer interface (BCI) has to be calibrated on every new user, and the cost is
higher still when the command has no natural counterpart, as with a robotic sixth finger.
We re-analysed public electroencephalography (EEG) recordings of \nSubjects{} stroke patients
who imagined grasping with the affected hand or flexing a sixth finger, and asked how much
calibration a decoder needs and whether recordings of other patients shorten it. The
dataset's own script fits its spatial filters on every trial before cross-validation: our
rebuild of that recipe scored \rqOneCspLdaReplicaPaperReplicaMean\,\% on hand against sixth
finger, and \rqOneCspLdaReplicaCvLeakfreeMean\,\% once each step was fitted on training
trials only. In pre-registered tests on \rqTwoHTwoAN{} patients, one re-centred Riemannian
transfer configuration, in which a single window of the patient counts no more than one of the thousands
contributed by other patients, reached \rqTwoRestAhSfRctTslrKOneMean\,\% on three classes after
a single trial of each movement, against \rqTwoRestAhSfTsLrKOneMean\,\% for the same
classifier trained from scratch on those four windows, with a larger area under the
learning curve ($p = \rqTwoHTwoAP$). It then stayed near \rqTwoRestAhSfRctTslrKAllMean\,\%,
and once every trial was used the patient-only decoder finished
\rqTwoRestAhSfTsLrMinusRctTslrKAllDiff{} points ahead, a gap we did not test. When the sixth finger was added as a new command, the Riemannian pair gave no clear gain before
therapy (the test was supported after therapy, so it did not replicate), and a decoder never
shown a sixth-finger trial detected one at least as well, so we claim no distinct
sixth-finger command. In exploratory
analyses, a decoder of sixth finger against rest scored \exFourTsLrCentralMean\,\% per
window but also fired on \exFourTsLrCentralFiresAhMean\,\% of hand-imagery windows, and in a
continuous replay it caught \exFiveFingerFingerOnMean\,\% of sixth-finger periods with
\exFiveFingerFalsePerMinMean{} false activations a minute. An ablation study gathers the
component comparisons the project had already run, each labelled pre-registered or
exploratory, and shows which choices mattered. All evaluation was offline.
